%% notes-tex/025-additive-baselines/sections/5-disjoint-transformer.tex

\section{The transformer on held-out screens\secstatus{todo}}
\label{sec:disjoint}

Four transformer runs have trained on arm Q. One is the 010 configuration. The other
three come from the joint-fitness-head branch (PR \#346, not yet on \texttt{main}):
they replace the learnable gene table with a fixed sequence embedding of matched
parameter count, read out through a perturbed CLS token, and train at a constant
learning rate for 30 epochs. They share the split and the graph penalty with GH 1640
and nothing else that a replicate would have to share, so they are four runs on one
split, not one run replicated.

%% SOURCE: experiments/025-solid-growth/scripts/additive_baselines_025_panels.py ->
%% tables/t3-disjoint-runs.tex, from the runs' W&B histories cached in
%% results/additive_baselines_025_arm_val_history.csv and
%% results/additive_baselines_025_disjoint_runs_history.csv.
\begin{table}[htbp]
\centering
\footnotesize
\caption{Every transformer run on arm Q. Epochs is the number logged: GH 1640 was
stopped by the 12 hour wall clock partway through epoch 36, and the IGB runs finished
their 30. The validation maximum is an upward-biased order statistic over the logged
epochs. No run has a test score. Configs of the IGB runs are on branch
\texttt{feat/025-fitness-joint-head}.}
\label{tab:disjointruns}
%% GENERATED by experiments/025-solid-growth/scripts/additive_baselines_025_panels.py
%% SOURCE: the result files named in that script's docstring. Do not edit by hand.
\begin{tabular}{llp{27mm}p{33mm}rrr}
\toprule
Run & Config & Gene input & Schedule, readout & Epochs & Val max (epoch) & Val last \\
\midrule
GH 1640 & \texttt{cgt\_s0\_q\_kl\_004} & learnable table & cosine, 010 schedule, mean & 36 & 0.199 (7) & 0.131 \\
IGB 2391132 & \texttt{cgt\_s0\_q\_kl\_emb\_017} & four-region sequence composite & constant $2.5 \times 10^{-4}$, perturbed CLS & 30 & 0.262 (2) & 0.222 \\
IGB 2391133 & \texttt{cgt\_s0\_q\_kl\_calm\_020} & CaLM codon embedding & constant $2.5 \times 10^{-4}$, perturbed CLS & 30 & 0.252 (7) & 0.208 \\
IGB 2391134 & \texttt{cgt\_s0\_q\_kl\_prot\_021} & ProtT5 protein embedding & constant $2.5 \times 10^{-4}$, perturbed CLS & 30 & 0.270 (3) & 0.135 \\
\bottomrule
\end{tabular}

\end{table}

%% SOURCE: experiments/025-solid-growth/scripts/additive_baselines_025_panels.py ->
%% additive_baselines_025_fig2_disjoint_runs.svg, from the same two history caches and
%% the arm Q val and test rows of results/additive_baselines_025.csv.
\tcfig{figures/additive_baselines_025_fig2_disjoint_runs.pdf}{Every transformer run
on the query-pair-disjoint split. \textbf{a}, validation Pearson per epoch on the
37{,}705 validation records of 43 held-out screens, with the additive ridge and the
three-seed MLP fit on the same validation part as horizontal lines; open circles
mark each run's best epoch. \textbf{b}, the same runs reduced to their best epoch
(open) and last logged epoch (filled). \textbf{c}, placeholder for the test
comparison that has not been run: the arm Q test nulls are drawn, and the empty
slots are the epoch 7 checkpoint of GH 1640 and three seeds of the 010
configuration, each to be scored on the 46 test screens at its best validation
epoch.}{fig:disjointruns}

\notesec{Every run peaks early and falls}
GH 1640 passes the additive validation null of 0.150 in epoch 4, peaks at 0.199 in
epoch 7, is last above the null in epoch 16 and ends at 0.131, below every null but
B3. The three IGB runs peak higher, 0.252 to 0.270 in epochs 2 to 7, and decline
from there. The composite and CaLM embeddings end at 0.222 and 0.208 and stay above
the ridge through their last five epochs; ProtT5 ends at 0.135, level with GH 1640.
Arm R has the opposite shape, a margin over its null that appears in epoch 2 and
holds for the rest of the run (Figure~\ref{fig:ladders}c).

\notesec{What the early peak does and does not say}
On validation, a transformer holds a margin over additivity for some epochs on
screens it has never seen, and the size of that margin depends on the gene input
and the schedule. Two things keep this from being a result. The maxima are order
statistics over 30 to 36 noisy epochs, so each is biased upward by an amount the
nulls, which are single fits, do not share. And validation is 43 screens; the test
part is a different 46, and Figure~\ref{fig:ladders}d shows how far one draw of
screens moves the null. A checkpoint scored on test at its validation-selected
epoch is the number that compares to Table~\ref{tab:heldout}, and no run has one.

\notesec{Hypothesis, untested}
The decline after the peak is consistent with the model fitting screen-specific
structure of the training pairs that does not transfer, the quantity B4 measures.
No measurement here separates that from ordinary overfitting of label noise, and
no mechanism is claimed.
